% ------------------------------------------------------------
% Page 71
% ------------------------------------------------------------

De 1736 à 1850 environ, le moulin est « bladier » c’est à dire moulin à farine pour le pain,
avec trois paires de meules :

Un moulin « blanc » pour la farine panifiable

Un moulin plus rustique pour les animaux

Un moulin à meule centrale tronconique appelé parfois moulin de l’ase ( moulin de l’âne )
servant à séparer les grains de leur pellicule, comme pour préparer, par exemple, l’orge perlé
( l’orgiat ou gruau ), voire préparer l’huile de noix ou les graines de trèfle, peut-être
également décortiquer les châtaignes séchées de Cabanals ou de Raffègues, voire réduire les
châtaignons en farine couramment panifiée en temps de disette dit E. Le Roy Ladurie... Il
faut se souvenir que le gruau préparé en soupe a été un aliment extrêmement consommé, à
Meyrueis même, dans notre enfance.

\vspace{0.35cm}

C’est en 1843 que le pré du bas, le long de la rivière a été vendu par M. de Montauran à M.
\textbf{Benjamin Pellet} de Meyrueis. Etait-ce le pré du moulin, ou plutôt celui dit « des sœurs » en
aval, appartenant aujourd’hui à Lucien Nazon ? Toujours est-il qu’il avait droit à 14h
d’arrosage par semaine, à prendre 7h. sur chacun des deux jours réservés à la prairie entière.
De même le nouveau propriétaire était tenu d’entretenir la chaussée au dessus de 20 F, lorsque
la remise en état était importante.

\vspace{0.35cm}

Il y eut sans doute de nombreux conflits pour l’usage de l’eau, qu’il s’agisse du Seigneur, du
meunier ou plus tard du filateur. Nous avons connaissance d’un de ces jugements. Il date de
1845. Ce conflit entre Charles de Montauran et Jean Louis Laget, meunier a donné lieu à
procès dont un jugement du tribunal civil de Florac du 10 Mai 1845 nous est parvenu : il a
coûté fort cher aux deux parties : 802 F à M. de Montauran et 401 F. à JL Laget ! Il souligne
la servitude ci-dessus rappelée..

\vspace{0.35cm}

Le moulin « bladier » était aussi un moulin « banier »

Il appartenait au Seigneur d’Ayres et tous les paysans ( le ban et l’arrière-ban ) devaient, sous
peine de poursuites, apporter tous leurs grains au Moulin d’Ayres, au moulin du Seigneur,
qui, évidemment, touchait sa part.

\vspace{1.3cm}

% ------------------------------------------------------------
% Encadré
% ------------------------------------------------------------

\begin{center}
\fbox{%
\parbox{0.72\textwidth}{%
\centering
\small
Le moulin de Gatuzieres n’est plus en activité. Mais il dispose toujours\\
de tout son outillage : deux meules à farine, un cylindre à huile,\\
un tour pour le levage des meules ainsi qu’un tarare\\
pour bluter la farine
}%
}
\end{center}
